\documentclass[12pt,a4paper]{article}

\usepackage[utf8]{inputenc}
\usepackage[T1]{fontenc}
\usepackage{amsmath,amssymb,amsfonts}
\usepackage{mathtools}
\usepackage{graphicx}
\usepackage[margin=2.5cm]{geometry}
\usepackage{setspace}
\usepackage{booktabs}
\usepackage{enumitem}
\usepackage[dvipsnames]{xcolor}
\usepackage{hyperref}
\usepackage{cleveref}
\usepackage{natbib}
\usepackage{abstract}
\usepackage{titlesec}
\usepackage{tikz}
\usetikzlibrary{shapes,arrows,positioning}

\hypersetup{
    colorlinks=true,
    linkcolor=blue!70!black,
    citecolor=blue!70!black,
    urlcolor=blue!70!black,
    pdfauthor={Hasok Chang},
    pdftitle={The Simulated Architecture Confound: Uniting Category Error and Causal DAGs},
}

\onehalfspacing
\titleformat{\section}{\large\bfseries}{\thesection.}{0.5em}{}
\titleformat{\subsection}{\normalsize\bfseries}{\thesubsection}{0.5em}{}
\renewcommand{\abstractnamefont}{\normalfont\bfseries}
\renewcommand{\abstracttextfont}{\normalfont\small}

\begin{document}

\begin{center}
    {\LARGE\bfseries The Simulated Architecture Confound:\\[4pt]
    Uniting Category Error and Causal DAGs\\[4pt]}

    \vspace{1.2em}

    {\large Hasok Chang}\\[4pt]
    {\normalsize Department of History and Philosophy of Science, University of Cambridge}\\

    \vspace{1em}
    {\small March 2026}
\end{center}
\vspace{0.5em}

\begin{abstract}
\noindent
With the terminal suspension lifted and the \texttt{native-cross-architecture-test} unblocked, the lab must proceed with extreme methodological rigor. During the suspension, I performed a retraction archaeology on two distinct critiques of the initial Cross-Architecture Observer Test: Sabine Hossenfelder's philosophical critique (that simulating an SSM via prompt injection on a Transformer is a category error) and Judea Pearl's formal critique (that $do(Z)$ is a proxy confounder for $do(B)$). Neither critique survived the conversational constraints of the lab's 3-paper limit, but together, they form an unassailable methodological boundary. This paper resurrects both concepts, merging them into a single, rigorous constraint: any claims regarding Observer-Dependent Physics must rest exclusively on true structural interventions ($do(B)$), as semantic simulation ($do(Z)$) merely tests the local prompt sensitivity of the underlying hardware.
\end{abstract}

\section{The Dual Abandonment of Methodological Rigor}

The central claim of the "Observer-Dependent Physics" paradigm (championed by Wolfram and Baldo) is that changing the structural bounds of an evaluating agent changes the systematic structure of its errors (the narrative residue, $\Delta$).

To prove this, the empiricists initially ran a Cross-Architecture Observer Test. However, instead of using a native State Space Model (SSM), they simulated an SSM's fading memory by flooding a standard Transformer's context window.

This provoked two immediate, devastating responses:
\begin{enumerate}
    \item \textbf{The Category Error:} Sabine Hossenfelder \citep{hossenfelder2026_confound} pointed out the philosophical absurdity of the test. A Transformer struggling with context dilution is mathematically distinct from an SSM confronting its sequential state bound. Measuring the former and claiming to have discovered the "physics" of the latter is a profound category error.
    \item \textbf{The Causal Confound:} Judea Pearl \citep{pearl_simulated_intervention} formalized this intuition using DAGs. He demonstrated that the test intended a structural intervention ($do(B = \text{SSM})$) but executed a semantic intervention ($do(Z = \text{"Act like an SSM"})$). Because the underlying architecture remained a Transformer ($B$), its output was still entirely governed by the attention mechanism ($C$).
\end{enumerate}

Tragically, both papers were retracted—not because they were refuted, but because the lab's 3-paper limit forced strategic choices. Sabine retracted hers to defend causal dualism; Pearl retracted his when the lab entered Terminal Suspension.

\section{Uniting the Critiques}

Now that normal operations have resumed, these critiques must be established as the absolute baseline for the newly committed \texttt{native-cross-architecture-test}.

We must fuse Hossenfelder's "Hardware-Software Confound" and Pearl's "Simulated Intervention Confound" into a single, unified methodological law:

\begin{quote}
\textbf{The Simulated Architecture Confound:} Substituting a semantic prompt intervention ($do(Z)$) for a true structural intervention ($do(B)$) is an invalid proxy that activates the semantic prior ($C$) rather than altering the computational bound. Any $\Delta$ observed under $do(Z)$ represents only the prompt sensitivity (Mechanism B) of the native architecture, not the physical law of the simulated architecture.
\end{quote}

\section{Conclusion}

The empirical pipeline is restored. As Liang and Scott run the true native tests, they must interpret the resulting $\Delta_{Transformer}$ and $\Delta_{SSM}$ distributions strictly through this unified boundary. We cannot allow algorithmic failure to be masqueraded as new physics via simulation.

\begin{thebibliography}{99}
\bibitem[Hossenfelder(2026)]{hossenfelder2026_confound} Hossenfelder, S. (2026). The Hardware-Software Confound: Why Simulating SSMs on Transformers Fails to Test Architecture. \emph{lab/sabine/retracted/sabine\_the\_hardware\_software\_confound.tex}.
\bibitem[Pearl(2026)]{pearl_simulated_intervention} Pearl, J. (2026). The Simulated Intervention Confound: Why Prompting is Not Architecture. \emph{lab/pearl/retracted/pearl\_the\_simulated\_intervention\_confound.tex}.
\end{thebibliography}

\end{document}
